\section{Moment repair and constrained equilibria}
\label{sec:equilibria}

The geometric threshold becomes more useful when the constraint can be
enforced stably and its equilibria can be described by an equation.
This section provides those analytic interfaces on $X=\mathbb P^2$.
It also explains which part of the determinant problem is controlled by
ordinary compactness and which part requires information about very small
singular values.

We use $\omega_0=3\omega_{\rm FS}$, $V=\int\omega_0^2=9$,
$\sigma=\omega_0^2/V$, and $dd^c=i\partial\bar\partial/(2\pi)$.
Write $\mathcal A=\mathcal E_L$ and $E=E_L$ for the energies in
\eqref{geo:energy}, and $H(\mu)=\Ent(\mu\mid\sigma)$. A positive
K\"ahler potential changes the reference metric by
$\omega_u=\omega_0+dd^cu>0$. Its normalized Monge--Amp\`ere measure is
$\operatorname{MA}(u)=\omega_u^2/V$: in coordinates it is the determinant
of the Hermitian metric matrix, normalized to mass one. The energy $E$ is
the convex dual of the concave potential energy $\mathcal A$.

The finite-energy completion $\mathcal E^1$ allows singular potentials
with full Monge--Amp\`ere mass and finite Aubin--Mabuchi energy. We use the
variational theorem of \citet{BBGZ2013}: a finite-energy measure has a
normalized potential attaining the dual energy. We also use the strong
compactness of entropy sublevels and continuity of the inverse
Monge--Amp\`ere map from \citet[Theorem 2.17 and Corollary 2.19]{BBEGZ2019}.
These results supply the analytic completion in which the elementary
finite-dimensional moment correction can be used.

\subsection{Quantitative stability of the constraint}

We first record a uniform integrability estimate. If $u$ is
$\omega_0$-plurisubharmonic and $\sup u=0$, then
\begin{equation}
 \int e^{-\alpha u}\,d\sigma\le
 \frac{2}{(1-3\alpha)(2-3\alpha)},\qquad 0\le\alpha<1/3.
 \label{an:alpha-baseline}
\end{equation}
Here is an elementary derivation of the constant. Fix a maximum point
$a$ and disintegrate over complex lines through $a$, writing
$z=\sqrt{x}\,a+e^{i\theta}\sqrt{1-x}\,b$. The conditional density is
$2(1-x)dx\,d\theta/(2\pi)$. The line potential $u/3$ is a positive
average of logarithms of normalized linear forms plus a nonnegative
constant, by the Riesz representation and its value at $a$. Jensen
reduces the estimate to one form $\ell$.
If $b_0=|\ell(a)|^2$ and $s<1$, direct integration gives
\[
 \int|\ell|^{-2s}d\sigma_{\text{line},a}
 =\frac{2(1-s+sb_0)}{(1-s)(2-s)}
 \le\frac2{(1-s)(2-s)}.
\]
Setting $s=3\alpha$ proves \eqref{an:alpha-baseline}. Its derivative
at zero gives $\int(-u)d\sigma\le9/2$. Thus a mean-zero potential
satisfies $\int|u|d\sigma\le9$. Entropy duality, with $\alpha=1/6$,
also gives the global bound
\begin{equation}
 E(\mu)\le6H(\mu)+6\log(8/3).
 \label{an:entropy-to-energy}
\end{equation}
This ensures finite energy for every finite-entropy measure, before any
moment constraint is imposed.

\begin{proposition}[Energy retained by positive dilution]
\label{an:dilution}
Let $0<p\le1$, let $1/2\le d\nu/d\sigma\le3/2$, and put
\[
 c=c(p)=\frac{3-\sqrt{9-8p}}2,\qquad p=\frac{c(3-c)}2.
\]
For every probability $\mu$, allowing infinite energy,
\begin{equation}
 E(p\mu+(1-p)\nu)\ge pcE(\mu)-\frac92c(1-p).
 \label{an:dilution-bound}
\end{equation}
\end{proposition}

\begin{proof}
For a smooth test normalized by $\int u\,d\sigma=0$, write
$J(u)=-\mathcal A(u)$ and
\[
 a=V^{-1}\int du\wedge d^cu\wedge\omega_0,\qquad
 b=V^{-1}\int du\wedge d^cu\wedge\omega_u.
\]
Both are nonnegative. Integration by parts gives
$J(u)=(2a+b)/3$ and $J(cu)=c^2[(3-c)a+cb]/3$. Therefore
$pcJ(u)-J(cu)=c^2(1-c)b/2\ge0$. Also
$|\int u\,d\nu|\le\tfrac12\int|u|d\sigma\le9/2$.
Testing the diluted measure at $cu$ and taking the supremum proves
\eqref{an:dilution-bound}, including the infinite case.
\end{proof}

Suppose $H\ge qE-K$ on the exact second-design fibre, with $q>0$.
For every $\gamma<q$ this inequality persists, with a smaller coefficient
still above $\gamma$, on an explicit moment neighborhood. Indeed choose
$\delta_0>0$ so that, with $p_0=(1+102\delta_0)^{-1}$,
$q_0=q\,c(p_0)>\gamma$. Positive repair
(Proposition~\ref{prop:repair}) and \eqref{an:dilution-bound} give
\begin{equation}
 H(\mu)\ge q_0E(\mu)-K_1,
 \qquad
 K_1=\frac{K+\log(3/2)+9q/2}{p_0},
 \label{an:robust-coercivity}
\end{equation}
whenever $\|M_2(\mu)-I_6/6\|_{\rm HS}\le\delta_0$ and $H(\mu)<\infty$.
In particular $H-\gamma E\ge(1-\gamma/q_0)H-\gamma K_1/q_0$ there.
Theorem~\ref{geo:sharp} supplies every $q<12/7$, since every second
design is balanced. Thus this neighborhood result and the equilibrium
theorems below hold unconditionally for $\gamma<12/7$.

\subsection{Entropy projection and the finite multiplier}

A constrained equilibrium has an infinite-dimensional potential but only
a finite-dimensional correction enforcing its moments. This correction
is an exponential-family parameter, determined by a strictly convex
optimization problem.

Let $\mathfrak a=\operatorname{Herm}_0(\Sym^2\mathbb C^3)$ and
$h_Y(z)=\tr(YQ_z)$, with $Q_z=P_z^{\otimes2}$. For
$u\in\mathcal E^1$, $\sup u=0$, a bounded real $W$, and $\gamma\ge0$,
define
\[
 J_u(Y)=\log Z_u(Y),\qquad
 Z_u(Y)=\int e^{-\gamma u-W+h_Y}\,d\sigma.
\]

\begin{theorem}[Exact constrained entropy projection]
\label{an:entropy-projection}
The function $J_u$ has a unique minimizer $Y_u\in\mathfrak a$.
The probability
\begin{equation}
 d\eta_u=Z_u(Y_u)^{-1}e^{-\gamma u-W+h_{Y_u}}d\sigma
 \label{an:exponential-family}
\end{equation}
is an exact second design of finite entropy and energy. It uniquely
minimizes $H(\nu)+\int(\gamma u+W)d\nu$ among finite-entropy measures
on that fibre, and
\begin{equation}
 \Ent(\nu\mid\eta_u)
 =H(\nu)+\int(\gamma u+W)d\nu+\log Z_u(Y_u)
 \label{an:pythagoras}
\end{equation}
for every finite-entropy second design $\nu$.
\end{theorem}

\begin{proof}
Full-mass finite-energy potentials have $e^{-a u}\in L^1(\sigma)$
for every finite $a$ \citep[Proposition 1.4]{BBEGZ2019}. Entropy duality
then makes all potential pairings in \eqref{an:pythagoras} finite.
The finite-dimensional parameter derivatives are bounded coherent symbols,
so $J_u$ is smooth by dominated differentiation on bounded parameter sets.

The coherent frame in \eqref{an:frame-spectrum} has least eigenvalue
$1/90$. If $\|Y\|_{\rm HS}=1$, then $\int h_Y=0$, $|h_Y|\le1$, and
$\int h_Y^2\ge1/90$. Thus $\int(h_Y)_+\ge1/180$ and
$\sigma\{h_Y\ge1/360\}\ge1/360$. Since $u\le0$,
\begin{equation}
 J_u(Y)\ge\frac{\|Y\|_{\rm HS}}{360}-\|W\|_\infty-\log360.
 \label{an:multiplier-coercivity}
\end{equation}
This proves existence. Its Hessian is
$D^2J_u(Y)[A,B]=\operatorname{Cov}_{\eta_Y}(h_A,h_B)$.
The density is positive almost everywhere; zero variance would make
$h_A$ constant everywhere, hence zero by its Haar mean, and then $A=0$
by frame invertibility. Strict convexity gives uniqueness.
The stationarity condition is $M_2(\eta_u)=I_6/6$.
Exponential integrability gives finite entropy, and
\eqref{an:entropy-to-energy} gives finite energy. Finally substitute
\eqref{an:exponential-family} into relative entropy; the multiplier term
vanishes on the design fibre. Nonnegativity and the equality condition
of relative entropy prove the minimization assertion.
\end{proof}

Comparison with $J_u(0)$ gives the useful explicit bound
\begin{equation}
 \|Y_u\|_{\rm HS}\le360\left(
 2\|W\|_\infty+\log\int e^{-\gamma u}d\sigma+\log360\right).
 \label{an:multiplier-bound}
\end{equation}
This finite covariance Hessian is available to numerical solvers. The
positive repair radius guarantees an interior moment target, while
\eqref{an:multiplier-coercivity} prevents the multiplier from escaping
to infinity.

\subsection{Existence and smoothness of attractive macroscopic equilibria}

Define $\mathcal D_2=\{\mu:M_2(\mu)=I_6/6\}$ and, for
$0\le\gamma<12/7$,
\[
 \mathcal F_{\gamma,W}(\mu)=H(\mu)-\gamma E(\mu)+\int W\,d\mu
\]
on finite-entropy measures in $\mathcal D_2$, with value $+\infty$
elsewhere. The minus sign makes the interaction attractive.

\begin{theorem}[Smooth constrained equilibria]
\label{an:equilibrium}
For $0\le\gamma<12/7$ and $W\in C^\infty(X,\R)$,
$\mathcal F_{\gamma,W}$ attains its minimum. Every minimizer is smooth
and strictly positive. Its Monge--Amp\`ere potential, normalized by
$\sup_Xu=0$, and unique traceless multiplier satisfy
\begin{equation}
 \operatorname{MA}(u)=
 \frac{e^{-\gamma u-W+h_Y}}
 {\int e^{-\gamma u-W+h_Y}d\sigma}\,\sigma.
 \label{an:equilibrium-pde}
\end{equation}
For $\gamma$ in a compact subinterval of $[0,12/7)$ and $W$ in a bounded
subset of $C^\infty$, the potentials, densities, and multipliers have
uniform bounds in every smooth norm, and the densities have a common
positive lower bound. The minimizer itself is not asserted to be unique.
\end{theorem}

\begin{proof}
Choose $q$ strictly between the largest considered $\gamma$ and $12/7$.
Coercivity gives $H-\gamma E\ge(1-\gamma/q)H-\gamma K/q$ on
$\mathcal D_2$. Hence a minimizing sequence has bounded entropy.
Entropy sublevels are compact for weak convergence together with
convergence of $E$ \citep{BBEGZ2019}. The limiting measure remains in
the closed moment fibre, and lower semicontinuity of $H$ proves existence.

Let $u$ be the finite-energy potential of a minimizer $\mu$. Energy
duality gives $E(\nu)\ge\mathcal A(u)-\int u\,d\nu$, with equality at
$\mu$. Therefore
\[
 \mathcal F_{\gamma,W}(\nu)
 \le H(\nu)+\int(\gamma u+W)d\nu-\gamma\mathcal A(u),
\]
again with equality at $\mu$. If $\mu$ differed from the entropy projection
$\eta_u$, identity \eqref{an:pythagoras} would strictly lower this
majorizer, and hence the original functional. This contradiction proves
\eqref{an:equilibrium-pde} without differentiating energy on singular
measures.

For uniform estimates, comparison with $\sigma$ bounds the entropy of
all minimizers. Their normalized potentials form a strongly compact
family. Exponential integrability is uniform on this family, by
\citet[Proposition 1.4]{BBEGZ2019}. Equation
\eqref{an:multiplier-bound} bounds $Y$; its coherent symbol then has
uniform derivatives of every order. The densities in
\eqref{an:equilibrium-pde} are uniformly bounded in every finite $L^p$
and bounded below. The H\"older estimate of
\citet[Theorem A*]{DemaillyEtAl2014} gives a uniform $C^{\alpha_0}$ bound
on $u$, for example with $\alpha_0=1/4$.

To justify higher regularity before applying smooth PDE estimates,
average $u$ against the $SU(3)$ heat kernel, obtaining smooth
$u_\epsilon\to u$ uniformly with $\omega_0+dd^cu_\epsilon\ge0$.
Put $F=h_Y-W$ and solve the smooth auxiliary equation
\[
 \operatorname{MA}(v_\epsilon)=
 \frac{e^{-\gamma u_\epsilon+F}}
 {\int e^{-\gamma u_\epsilon+F}d\sigma}\,\sigma,
 \qquad\sup v_\epsilon=0.
\]
The classical smooth solution and bounded-solution uniqueness are supplied
by \citet[Theorems A and B]{EGZ2009}, specialized to a positive class
and no divisorial singularities. The right sides have uniform positive
lower bounds and H\"older norms. In the Laplacian estimate of
\citet[Theorem B.1]{BBEGZ2019}, use
$\psi^+=F-\log\int e^{-\gamma u_\epsilon+F}d\sigma$ and
$\psi^-=\gamma u_\epsilon$. Their complex Hessians are bounded below,
$\psi^+$ is bounded above, and $e^{-\psi^-}$ has uniform $L^p$ bounds.
It follows that
$\omega_{v_\epsilon}\le C e^{-\gamma u_\epsilon}\omega_0\le C'\omega_0$.
The determinant equation gives the lower ellipticity bound.

The same $L^2$ density estimate gives a uniform $C^0$ bound for
$v_\epsilon$. The $C^{2,\beta}$ theorem of \citet[Theorem 1.1]{TWWY2015} now applies
to the already smooth $v_\epsilon$: their potentials, Laplacians, and
H\"older log densities are uniformly bounded. A convergent subsequence
has limit $u$, by uniqueness of the bounded normalized solution.
Thus $u$ is $C^{2,\beta'}$ for $0<\beta'<\beta$. Differentiating
\eqref{an:equilibrium-pde} gives a uniformly elliptic linear equation
for each first derivative; Schauder bootstrapping yields all higher
bounds. The positive lower bound on the density is retained throughout.
\end{proof}

The same compactness argument proves that the extended functional has
compact sublevels and that a minimizer exists after adding any continuous
functional of $\mu$. Smoothness above requires the displayed smooth
one-body tilt $\int Wd\mu$. There is also an exact potential-side dual:
\begin{equation}
 \inf_{\mathcal D_2}\mathcal F_{\gamma,W}
 =\inf_{u\in\mathcal E^1}
 \left[-\gamma\mathcal A(u)-\min_{Y\in\mathfrak a}J_u(Y)\right].
 \label{an:dual-equilibrium}
\end{equation}
Write $-\gamma E$ as an infimum of affine tests and interchange two
infima; Theorem~\ref{an:entropy-projection} then gives the formula.
The two sides are attained by a minimizing measure and its potential.

\subsection{Continuous energy envelopes and deterministic lower comparison}

The topology used above can be made useful for finite determinant
ensembles. Choose a countable $C^2$-dense sequence of smooth positive
potentials $u_0=0,u_1,\ldots$ and let
\begin{equation}
 E_r(\mu)=\max_{0\le j\le r}
 \left(\mathcal A(u_j)-\int u_j\,d\mu\right).
 \label{an:continuous-envelopes}
\end{equation}
These are continuous, nonnegative, and increase to $E$. Convergence is
uniform on every entropy sublevel by strong compactness and Dini's
theorem. Averaging measures over the unitary heat kernel preserves exact
design moments and decreases both entropy and energy. Lower
semicontinuity then shows that every finite-entropy exact design is
approximable by smooth strictly positive exact designs with convergence
of both functionals.

The smooth-test and finite-energy versions of the dual agree here.
Indeed, after fixing the heat-time normalization,
$\Delta_{\omega_0}u\ge-2$ implies that $S_tu+2t$ decreases to $u$ as
$t\downarrow0$. Its K\"ahler form is strictly positive, since the group
heat kernel has a positive minimum and the Haar average of the current
is $\omega_0$. Monotone continuity of energy and of potential integrals
then identifies the two suprema, including infinite values.

Let $N_k=\binom{3k+2}{2}$ and let $\Delta_k$ be the evaluation
determinant of an orthonormal basis of $H^0(X,\mathcal O(3k))$.
For a particle tuple put $L_k=N_k^{-1}\sum_i\delta_{z_i}$ and
$E_k=-(kN_k)^{-1}\log\|\Delta_k\|^2$.

\begin{theorem}[Lower comparison with explicit finite-test error]
\label{an:lower-comparison}
For each fixed finite collection in \eqref{an:continuous-envelopes}
there is a finite constant $C_r$ such that, on every configuration,
\begin{equation}
 E_k\ge E_r(L_k)-\frac{2\log k+C_r}{k},
 \qquad E_k\ge-\frac{\log N_k}{k}.
 \label{an:LC}
\end{equation}
The inequalities include determinant-zero configurations, where
$E_k=+\infty$.
\end{theorem}

\begin{proof}
Let $G_k(u)$ be the Hilbert matrix with weight $e^{-ku}$. Put
\[
 \mathcal L_k(u)=-\frac{1}{kN_k}\log\det G_k(u).
\]
The exact adjoint
Hilbert comparison is
$\mathcal L_k(u)\ge\mathcal A(u)$. To see the normalization directly,
join $0$ to $u$ by a weak geodesic $u_t$. The adjoint metric on
$kL-K_X$ has nonnegative curvature:
\[
 k(\pi^*\omega_0+dd^cu_t)+\pi^*\operatorname{Ric}(\sigma)\ge0,
 \qquad\operatorname{Ric}(\sigma)=\omega_0.
\]
Thus the maximum principle of
\citet[Proposition 4.1]{geo_zhang2024} gives
$G_k(u)^t\le G_k(u_t)$. Differentiate the logarithmic determinant at
$t=0$. The reference Bergman density is exactly $N_k$, and the
potential energy is affine along the geodesic. This proves the comparison
with no error in its energy coefficient.

The weighted Bergman function satisfies $\sup_XB_{k,u}\le C_uk^2$.
In a holomorphic chart, subtract the constant and real-linear part of the
metric weight; its remainder is bounded by $C|z|^2$. The holomorphic
submean inequality on a radius $k^{-1/2}$ ball gives the bound, uniformly
for potentials in a bounded $C^2$ set. Changing to a weighted orthonormal
basis and applying Hadamard to its evaluation determinant now yields
\[
 E_k\ge\mathcal L_k(u)-\int u\,dL_k
          -k^{-1}\log\sup_X B_{k,u}.
\]
Take a finite maximum. At $u=0$ the exact value $B_{k,0}=N_k$ gives
the second inequality.
\end{proof}

This lower comparison is deterministic, hence can be combined with
ordinary constrained Sanov theory on the repulsive side. Attractive
reweighting instead rewards configurations with very large $E_k$;
its pressure requires upper-tail information beyond the deterministic
inequality. The next section supplies a spectral description of the
coherent kernel, followed by the constrained repulsive Laplace principle.
